\section*{ID6776: Next Steps I Can Run for You Immediately (Pick Any): a\textasciicircum{}rm obs\_m}
\paragraph{Metadata.}
PiID: 5544959778377 | RTEID: RTE:P26312.7128:T5.6561 | UUID: 588adbb5-4f77-5800-9df2-ba03b757a596

\paragraph{Canonical Statement.}
\#\# Next steps I can run for you immediately (pick any) 1. Run the same analysis specifically on your **starburst photos** (if those images become available in the runtime) to extract clean m-peaks and show the doubling relation with spherical harmonics. 2. Produce a **deconvolution attempt**: assume a simple PSF (e.g., pupil with N vanes) and try to invert \textbackslash{}(a\textasciicircum{}\{\textbackslash{}rm obs\}\_m = T\_m a\textasciicircum{}\{\textbackslash{}rm true\}\_m\textbackslash{}) to estimate underlying TZP spherical-harmonic components for the 3D render. (This needs careful regularization.) 3. Build a small GUI-like notebook cell that accepts an uploaded image and runs the pipeline (so you can run locally when you can upload more images). I can paste that as a script you can run. 4. Run a chromatic-phase analysis to more carefully test whether the purple seam / redshift you observed is optical chromatic dispersion or structural (this requires high-res RGB separation and possibly calibration).

\paragraph{Canonical Equation.}
\[
a^{\rm obs}_m = T_m a^{\rm true}_m
\]

\paragraph{Dictionary.}
Relation: a\textasciicircum{}rm obs\_m = T\_m a\textasciicircum{}rm true\_m [ID6776]

\paragraph{Encyclopedia.}
Next Steps I Can Run for You Immediately (Pick Any): a\textasciicircum{}rm obs\_m is an algebraic definition, written a\textasciicircum{}rm obs\_m = T\_m a\textasciicircum{}rm true\_m. It is an algebraic definition expressing the quantity directly in terms of the variables on its right-hand side. It is built from T\_m, a, r, u, m, defining a\textasciicircum{}rm obs\_m. Within the Trawin Zero Point registry it serves as one element of the framework's mathematical narrative.
